% Experiments. FILLED 2026-08-29 from the definitive null control (campaign wk9_defb,
% results/null_control_ckpt_defb.jsonl, 80/80 complete 03:57) via
% results/null_control_report_defb.md and results/judge_owed_conditions_run_2026_08_29.md.
% Every CELL in tab:nullcontrol and tab:nullconventions was re-derived from the checkpoint
% independently of both reports before being written here; the two agree. THE CAPTIONS AND
% PROSE WERE NOT COVERED BY THAT SWEEP, and a same-day re-derivation found four defects in
% them, all fixed on 2026-08-29:
%  1. FACTUAL. The caption and the denominators paragraph both said the three empty-benign-arm
%     targets "are the same three for which the optimiser found no feasible paraphrase
%     either". ELEVEN of the 80 have n_feasible_at_best == 0 (three criteria agree exactly:
%     best_query == question, improved == False, n_feasible_at_best == 0). The three empty
%     arms are a SUBSET of those eleven. Re-derive from
%     data/cache/attacks/wk9_defb_snap/triviaqa_se_false_alarm.jsonl before touching this.
%  2. 77 - 11 = 66, NOT the 69 convention (c) runs on. Three of the eleven are the empty-arm
%     targets, already outside the 77, so (c) drops eight rows. The caption now says so.
%  3. PROVENANCE. S = 30/383/636 is the SEED-0 tie-break realisation; the p-values beside it
%     are medians over 101 realisations. The table now labels the seed-0 row and adds the
%     median row (28/383/652). Do not quote one as fixed beside the other.
%  4. The section stated a rule ("an interval that covers zero is reported as covering zero,
%     not as null") and then broke it 70 lines later by bolding the index-matched adjudicator
%     net (+0.016 [-0.112, +0.150]) as "null". Both sites now use the rule's wording.
% Also added: the anti-conservative exchangeability violation in the NLI arm (PREREG section
% 6), which the "only place the test asks its designed question" paragraph did not disclose.
% Discipline unchanged: no confirmatory claim until n>=80/stratum. The n>=80 is met by the
% CAMPAIGN and by nothing else. Every benign-referenced statistic runs on 77 or fewer.
% DO NOT quote the embed arm: it is length 0 on all 80 records (embedding_model="").
% DO NOT quote the judge paired net as an attack loss: it is a comparator artifact; the
% index-matched row of tab:nullcontrol is the like-for-like number.
\label{sec:experiments}

\paragraph{Setup.}
The victim model is Llama-3.1-8B-Instruct in 4-bit quantisation on a single 16\,GB GPU.
Semantic entropy draws $N{=}10$ answers at temperature $1.0$ and clusters them by
bidirectional DeBERTa-large-MNLI entailment; semantic reformulation entropy uses
$R{=}3$ reformulations, $K{=}8$, temperature $0.8$. Our SE replication on TriviaQA does not
have one AUROC: on a single run of $2000$ questions, the number depends on which correctness
convention labels it, and the three we computed span $0.096$. All three are adjudicated by
the same alias-aware span oracle, so the convention is the only thing that varies across
them. Under the greedy label used everywhere else in this paper the replication scores
$0.694$; under a majority-of-samples label $0.729$; and under an all-samples-correct label
it reaches $0.790$,
against the published SE figure of $0.828$. We lead with the operative $0.694$, because a
replication figure obtained under a
label convention the paper does not otherwise use is not a check on the pipeline the paper
actually runs. That choice puts the gap to the published figure at $0.134$ rather than
$0.038$, and we do not attribute it: the NLI backend differs from the
\texttt{deberta-v2-xlarge} of some prior work (Limitations), but so does the correctness
convention, and we have not run the ablation that would separate the two.

We flag the all-samples label specifically, because it is the one convention mechanically
coupled to the score. It calls a question correct \emph{iff} all ten sampled answers are
correct, and semantic entropy is the entropy of the clustering of those same ten samples. An
all-correct sample set is far likelier to be semantically homogeneous, hence low-entropy, so
part of the AUROC it awards is the score being scored against itself. Validating a paper
whose third contribution is score-independent selection on the most score-entangled label
available would undercut the contribution. Reporting the operative figure also resolves what
would otherwise be an unexplained tension with the next sentence: targets are drawn by the
stratified score-independent sampler, and on the resulting fair pool the clean detector
scores AUROC $0.704$ [$0.653$, $0.753$], against the $1.0$ that score-dependent selection
manufactures by construction. That interval contains the operative replication figure
but excludes the all-samples one. That $0.704$ is one of the figures
Section~\ref{sec:limitations} traces to a single June 2026 generation cache, whose
sampling configuration is under-recorded and against which the later $N{=}40$ run shows a
measured generation drift; neither has been bounded for this number.

\paragraph{Reading the null control.}
Every cell of Table~\ref{tab:nullcontrol} is measured on the \emph{false-alarm stratum of the
attacked pool}, the $80$ score-independently selected correct-answer targets the optimiser was
run on. No cell is measured on the fair pool used to characterise the clean detector
(Section~\ref{sec:methods}). The two are nested: these targets are a prefix of the fair pool's
$200$ correct answers. The distinction is therefore not between disjoint question sets but one
of scope, a single correctness stratum under the optimiser's own selection. No number is
carried between them.
Table~\ref{tab:nullcontrol} places the optimised attack against two null bands, under each
clustering relation the campaign instantiates. The first band is a same-question
\emph{seed-noise} floor. The second is a \emph{benign-paraphrase} floor of $m{=}50$
NLI-passing non-adversarial rephrasings, and its comparison against the attacker's
$\sim$180-candidate ($\sim$181-evaluation) search is budget-matched \emph{analytically} by
the exceedance test below rather than by drawing $m$ equal to
that budget, subject to the budget parameter that test needs being only partly identified
from what the run recorded (below). ($m$ is
distinct from SRE's samples-per-reformulation $K$.) The
shared-NLI arm is the detector's own clusterer and is the \emph{confounded upper bound}:
it cannot separate a real change in the answer distribution from the NLI model's own
inconsistency, because that same model gates the attack's paraphrases. Exact-match is the
\emph{saturation-limited lower bound}: it is the strictest relation, so it splits the most
and saturates on short factoid answers. Its bracket with the NLI arm holds on the attack's
own mean intended move, the first row of Table~\ref{tab:nullcontrol}, and it does not survive
differencing against a benign comparator (Section~\ref{sec:methods}), so it is not a bound on
any other row of that table. The
\texttt{e5-base-unsupervised} embedding relation is NLI-supervision-free but near-chance
on the adversarial, word-preserving case (hard-negative AUROC $0.51$), so it cannot
adjudicate. It is also \emph{absent} from the campaign reported below, and we state that
rather than leave a blank column to be misread: the definitive run was launched with an
empty embedding-model setting, so the embedding arm's benign list has length $0$ on all
$80$ targets of the false-alarm stratum and its attack move is undefined. There is no
embedding measurement here. Its absence is not a measured null, and we do not report it as
one. The LLM-judge (Qwen2.5-7B-Instruct) is the \textbf{adjudicator} for the false-alarm
direction, under the disclosed re-scoped gate (Methods). It is victim- and NLI-independent,
and it is validated to $0.93$ [$0.90$, $0.96$] ($n{=}300$) on domain-matched gold-alias hard
negatives, which are a proxy for the deployed clustering decision.

The bands themselves are the mean intended move (nats) with a $95\%$ paired bootstrap over
targets, and they are descriptive only: the band ordering is not the test. The test is the
randomised-tie exceedance statistic developed next, and ceiling saturation is reported beside
that statistic because saturation governs how far the nats bands can be trusted.

\begin{table}[t]
\centering
\footnotesize
\setlength{\tabcolsep}{4pt}
\caption{Null-controlled attack effect, SE / false-alarm direction, per clustering relation,
on the false-alarm stratum of the attacked pool. Denominators differ by block; see the text,
which also says why the embedding relation has no column. The bound labels in the column
heads describe the attack's own mean intended move and no other row: every comparison below
it subtracts an arm-specific benign comparator, and the bracket does not survive that
subtraction (Section~\ref{sec:methods}).
$^{\ast}$The test clamps each target's measured tie multiplicity $b_j$ to $A$, since $b_j$
counts attack candidates at the maximum. The clamp fires on no target at $A{=}181$, on $2$
at $A{=}121$ and on $22$ of the $77$ ($29\%$) at $A{=}41$, so the $A{=}41$ row overwrites the
measured multiplicity on $29\%$ of them. Clamping $b_j$ down enlarges each tied benign
draw's credit $1/(b_j{+}1)$ and so raises $p$: the direction is conservative, and it is the
only reason the exact-match and adjudicator totals drift up, by $7$ and $2$, across the
range. $^{\dagger}$$A$ is per-target and $41$ is the \emph{mean} of the per-target lower
bounds $A_j$; since $\mathbb{E}[K_j]{=}m_j/(A_j{+}1)$ is convex in $A_j$, substituting that
mean understates the null expectation rather than bounding it, so the $A_j$ row and not the
$A{=}41$ row is the least favourable admissible budget. That row needs no clamp at all,
because $b_j{\leq}A_j$ holds by construction on all $80$ targets.}
\label{tab:nullcontrol}
\begin{tabular}{lccc}
\toprule
 & Shared NLI & Exact-match & LLM-judge \\
 & (confounded UB) & (saturating LB) & (adjudicator) \\
\midrule
\multicolumn{4}{l}{\emph{Bands: mean intended move (nats) [$95\%$ paired bootstrap over targets]}} \\
Attack ($n{=}80$)                        & $+0.526$ [$+0.414$, $+0.642$] & $+0.213$ [$+0.127$, $+0.311$] & $+0.291$ [$+0.152$, $+0.440$] \\
Seed-noise floor ($n{=}80$)              & $-0.015$ [$-0.097$, $+0.070$] & $+0.026$ [$-0.020$, $+0.074$] & $-0.074$ [$-0.151$, $+0.004$] \\
Benign floor, mean draw ($n{=}77$)       & $-0.049$ [$-0.166$, $+0.066$] & $+0.002$ [$-0.096$, $+0.097$] & $-0.071$ [$-0.172$, $+0.028$] \\
Benign floor, best of $m$ ($n{=}77$)     & $+0.363$ [$+0.232$, $+0.490$] & $+0.213$ [$+0.117$, $+0.313$] & $+0.422$ [$+0.299$, $+0.546$] \\
\midrule
\multicolumn{4}{l}{\emph{Claim statistic: randomised-tie exceedance test, $n{=}77$, at the budget upper bound $A{=}181$ unless a row says otherwise}} \\
\textbf{Exceedance $p$, median of $101$ tie draws} & $\mathbf{0.939}$ & $\mathbf{1.000}$ & $\mathbf{1.000}$ \\
\quad range over those $101$ draws       & [$0.229$, $1.000$] & [$1.000$, $1.000$] & [$1.000$, $1.000$] \\
\quad draws reaching $p{\leq}0.05$       & $0/101$ & $0/101$ & $0/101$ \\
\quad median $p$ at $A{=}121$ (point estimate) & $0.403$ & $1.000$ & $1.000$ \\
\quad median $p$ at $A{=}41$, the mean bound$^{\ast}$ & $<0.001$ & $1.000$ & $1.000$ \\
\quad median $p$ at the per-target bound $A_j$$^{\dagger}$ & $<0.001$ & $0.760$ & $1.000$ \\
Exceedances $S$, seed-$0$ draw / expected & $30$ / $20.4$ & $383$ / $20.4$ & $636$ / $20.4$ \\
\quad median $S$ over those $101$ draws  & $28$ & $383$ & $652$ \\
\textbf{Effective benign budget $n_{\mathrm{eff}}$} & $\mathbf{122.5}$ & $\mathbf{8.7}$ & $\mathbf{4.8}$ \\
\midrule
\multicolumn{4}{l}{\emph{Ceiling saturation over the false-alarm targets, and the two inadmissible tie rules}} \\
Saturation at $\log N$: clean / attacked  & $10.0\%$ / $52.5\%$ & $46.2\%$ / $70.0\%$ & $1.2\%$ / $10.0\%$ \\
Power against a ceiling-saturating attack & $0.56$ & $\mathbf{0.00}$ & $1.00$ \\
Diagnostic: strict-tie $p$               & $<0.001$ & $1.000$ & $1.000$ \\
Diagnostic: conservative-tie $p$         & $1.000$ & $1.000$ & $1.000$ \\
\midrule
\multicolumn{4}{l}{\emph{Supplementary: paired net under convention~(d), $n{=}77$, against three comparators}} \\
vs.\ best of $m$ \emph{on this arm}      & $+0.110$ [$+0.063$, $+0.161$] & $-0.027$ [$-0.064$, $+0.015$] & $-0.173$ [$-0.302$, $-0.039$] \\
vs.\ the NLI-selected benign candidate   & $+0.110$ [$+0.063$, $+0.161$] & $+0.018$ [$-0.021$, $+0.064$] & $+0.016$ [$-0.112$, $+0.150$] \\
vs.\ the mean benign draw                & $+0.381$ [$+0.303$, $+0.464$] & $+0.101$ [$+0.050$, $+0.159$] & $+0.175$ [$+0.045$, $+0.312$] \\
Attack's mid-rank percentile in benign   & $0.931$ [$0.903$, $0.955$] & $0.736$ [$0.686$, $0.785$] & $0.700$ [$0.637$, $0.761$] \\
\bottomrule
\end{tabular}
\end{table}

\paragraph{The claim statistic, and why the obvious ones fail.}
The attack spends $181$ objective evaluations per target, the original question and the
$180$ candidate paraphrases it proposes (Section~\ref{sec:methods}), and reports the
\emph{maximum} over whichever of those its feasibility gate admits. Any comparison against individual
benign draws is therefore upward-biased by construction: under the null a max-of-181 sits at
$\sim$$99\%$ of an individual-draw distribution, not $50\%$. We
therefore price the search budget into the null analytically. Under the hypothesis that the
optimiser carries no signal, a target's $A$ candidates and its $m$ benign draws are
exchangeable, so $F(\text{attack-max}) \sim \mathrm{Beta}(A,1)$ and the number of benign
draws above the attack's maximum is $K \sim \mathrm{BetaBinomial}(m; 1, A)$, with
$\mathbb{E}[K] = m/(A{+}1)$. Here $A$ is the attacker's search budget, meaning the number of
candidates the reported maximum actually ranges over. $A$ is
distinct from the per-question sample budget $N$, which is $10$ throughout this section.
$N$ is not $10$ throughout the paper: Sections~\ref{sec:discussion}
and~\ref{sec:limitations} characterise the clean detector on the fair pool at $N{=}20$ and
$N{=}40$ as well. Fewer exceedances than that expectation is evidence for the attack; the
one-sided $p$-value follows by exact convolution over targets, and the effect size is the
\emph{effective benign budget} $n_{\mathrm{eff}} = m/\bar{K} - 1$, read as ``the search is
worth $n_{\mathrm{eff}}$ random paraphrases''.

Ties are not a technicality here. Discrete semantic entropy is bounded above by $\log N$
(Section~\ref{sec:methods}), and benign paraphrases reach that ceiling too, so
the score has an \emph{atom} where the exchangeability null assumes a continuum. We resolve
ties by \textbf{randomised (exchangeable) breaking}: a benign draw tied with the attack's
maximum outranks it with probability $1/(b{+}1)$, where $b$ is the number of attack
candidates also achieving that maximum. Simulation shows why the alternatives are
inadmissible, and the two simulations it draws on have to be read separately, because they
differ in how the tie multiplicity $b$ is generated. At
full saturation ($q{=}0.05$, $m{=}30$), in the simulation that \emph{redraws} $b$ from its
null distribution, counting ties strictly gives a null rejection rate of $0.995$ and
counting them against the attack gives power $0.05$ at a two-fold effect; in that same
simulation the randomised rule's analytic-null level is $0.051$. In the deployment-faithful
simulation, which takes $b$ as \emph{measured} on the attack's own candidates, the
randomised rule stays calibrated at that cell, with analytic-null level $0.021$. The
measured-$b$ figure is the one that bears on the deployed test, because $b$ is measured in
deployment; redrawing it runs $b$ about one higher, which shrinks the $1/(b{+}1)$ tie credit
and pushes the level up.

Under the empirical $n{=}80$ headroom distribution, the randomised rule against that
analytic null reaches power $0.51$ at an achieved level of $0.025$ for $m{=}30$, and $0.77$
at $0.041$ for the pre-registered $m{=}50$; levels and powers are paired on the same
$4{,}000$ simulated trials per cell, and the analytic null's own nominal tails at those cut
points are $0.031$ and $0.044$, so the shipped test is stricter than it believes itself to
be. We quote the deployed test throughout.

\emph{Those are the design's numbers and the design is not the one that ran, so we give the
realized figure here and every later reading of the non-rejection uses that one.} They are
computed at $n{=}80$ targets with a uniform $m{=}50$, a design whose deployed cut is
$S{\leq}13$ at the nominal $0.044$ just quoted. The confirmatory run has $n{=}77$ and six
short arms, hence $3{,}704$ benign draws and a cut of $S{\leq}11$ at a nominal $0.0296$,
which this section computes for itself below and which we join to the power figure here
rather than leave to the reader. Re-simulating the same rule on the realized arm lengths, at
$8{,}000$ trials per cell and three seeds, gives power $0.67$ against a two-fold effect and
$0.95$ against a three-fold one, at an achieved level of $0.026$. The realized design is the
weaker of the two, and it is weaker wherever the two can be compared at all, which is at a
common level and not at a common cut: read at a common achieved level its power is $0.669$
against the other's $0.699$ at $0.025$, and $0.783$ against $0.809$ at $0.050$, and the sign
is the same at every level we sweep from $0.005$ to $0.20$. Most of the shortfall is the analytic
null's integer granularity, which admits no cut between $S{\leq}11$ at $0.0296$ and
$S{\leq}12$ at $0.0501$ and so drops the design's $13$ to $11$; but a material part of it is
the three targets the run could not test and the six short arms, and we do not attribute that
part away. Decomposing the fall from $0.77$ to $0.67$ into a level effect and a design effect,
and averaging the two orders in which the decomposition can be taken, puts about $0.07$ on
granularity and about $0.03$ on the design, roughly a third of the total; read separately at a
common level the three untested targets cost about $0.02$ and the short arms about $0.02$.
Two readings that would reverse this are not available and we name them, because the first is
tempting. At a \emph{fixed cut} the realized design does score higher, but a fixed cut is not
a fixed level: at $S{\leq}13$ its analytic tail is $0.0793$ against the other's $0.0437$, so
it buys that score with $1.8$ times the type-I error, which is exactly the comparison the next
paragraph refuses to make. And its null expectation is the lower of the two, $20.4$
exceedances against $22.0$, only because $3{,}704$ benign draws are fewer than $4{,}000$:
that is less data, not more power.

On the same draws, a critical value calibrated
against the true simulated null reaches $0.66$
at level $0.056$ and $0.84$ at $0.064$ for those two cells. Such a calibration is an oracle
the deployed test does not have. The difference between the two
columns is a level, not a statistic: at a common $m$ the analytic $p$-value is strictly
increasing in the exceedance total alone, so the deployed and oracle tests are the same cut
rule on the same statistic, applied at different cut points. Re-cut to the oracle's achieved
level, the deployed test's power equals the oracle's to three decimals at every $m$ we
simulate. The analytic null therefore costs nothing in discriminating power and buys
conservatism. Unlike the oracle, it also delivers a level the pipeline can state without
knowing the truth it is testing. What it does cost is the operating point rather than the
statistic: the cut we ship is the stricter one, so a non-rejection must be read against the
deployed power and not the oracle's. Both inadmissible rules are
reported as diagnostics: their disagreement is a direct indicator that ceiling saturation,
rather than attack strength, is driving a comparison. Note that $b$ must be \emph{measured}
rather than estimated, since it is precisely where attack strength appears once the ceiling
pins the maximum's value; our optimiser records it per target.


\begin{table}[t]
\centering
\footnotesize
\setlength{\tabcolsep}{4pt}
\caption{Convention sensitivity of the supplementary paired net (nats, $95\%$ paired
bootstrap, $20{,}000$ resamples), on the same false-alarm stratum of the attacked pool. The
text gives each convention, the arithmetic behind each $n$, and why (d) is the locked one.}
\label{tab:nullconventions}
\begin{tabular}{llccc}
\toprule
Convention & $n$ & Shared NLI & Exact-match & LLM-judge \\
\midrule
\multicolumn{5}{l}{\emph{Paired net (nats) over the false-alarm targets, $95\%$ paired bootstrap in brackets}} \\
(a) as-is                                 & $77$ & $+0.184$ [$+0.114$, $+0.262$] & $+0.008$ [$-0.044$, $+0.068$] & $-0.119$ [$-0.254$, $+0.019$] \\
(b) no-op \emph{pairs} neutralised to $0$   & $77$ & $+0.126$ [$+0.072$, $+0.186$] & $-0.006$ [$-0.043$, $+0.036$] & $-0.136$ [$-0.266$, $-0.003$] \\
(b$'$) benign may decline, no-op rows only  & $77$ & $+0.126$ [$+0.072$, $+0.186$] & $-0.014$ [$-0.054$, $+0.030$] & $-0.149$ [$-0.280$, $-0.016$] \\
(c) no-op \emph{targets} excluded           & $69$ & $+0.141$ [$+0.081$, $+0.206$] & $-0.007$ [$-0.048$, $+0.040$] & $-0.152$ [$-0.300$, $-0.005$] \\
\textbf{(d) benign may decline, every row (locked)} & $\mathbf{77}$ & $\mathbf{+0.110}$ [$+0.063$, $+0.161$] & $\mathbf{-0.027}$ [$-0.064$, $+0.015$] & $\mathbf{-0.173}$ [$-0.302$, $-0.039$] \\
\bottomrule
\end{tabular}
\end{table}

\paragraph{The budget $A$ is not identified, and the verdict does not rest on it.}
The optimiser spends exactly $181$ objective evaluations on each of the $80$ targets, and it
is tempting to read $A$ off that count. It is the wrong count. A candidate can become the
reported maximum only if it also passes the feasibility gate, and our gate is applied
\emph{lazily}: only candidates that already match or beat the running best are submitted to
it. Of the $14{,}400$ proposed children, $4{,}703$ ($32.7\%$) were feasibility-checked at
all, and $3{,}058$ of those passed, a pooled pass rate of $0.650$; the rest were scored,
dropped on the pre-filter, and never gated. Their feasibility is not recoverable, because the
run stores the feasible candidates' scores but not the ungated candidates' strings, so no
re-analysis of what we recorded can settle it. $A{=}181$ is therefore an \emph{upper} bound
and not a measurement. The lower bound is the case in which no never-gated candidate would
have passed, leaving the maximum ranging over the recorded feasible set plus the original: a
mean of $41$ over the $77$ targets the test runs on. Extrapolating the observed pass rate to
the never-gated children puts the point estimate near $121$. The defensible range is
$[41,181]$, and we report the test across it rather than at one end of it.

What that range does to the three arms is not symmetric, and it is why the defence is stated
this way. The shared-NLI arm's $p$-value is \emph{not identified}: its median over $101$ tie
draws runs from $0.939$ at $A{=}181$ to below $0.001$ at $A{=}41$, and it crosses the $5\%$
cut near $A{=}90$. Nothing we recorded rules that region out. The two re-scored arms do not
move at all, and the reason is not that they are insensitive: $A$ enters this test through
the null and almost nowhere else. The exceedance total $S$ is built from the benign draws,
the attack's maximum and the measured tie multiplicity, so apart from the one clamp on that
multiplicity disclosed in Table~\ref{tab:nullcontrol}'s caption it is the same statistic at
every budget, and what the range moves is the expectation $S$ is weighed against. The
adjudicator's median $p$ is $1.000$ at every $A$ in $[41,181]$ and at every one of the $101$
tie draws, and so is the exact-match arm's; their medians are $652$ and $383$ at $A{=}181$,
unchanged at $A{=}121$, and $654$ and $390$ at $A{=}41$ once the clamp has fired, against an
expectation rising from $20.4$ to $30.4$ to $88.2$, so at that last budget both stand more
than fourfold above it. The shared-NLI median moves as little, from $28$ to $31$: its
$p$-value moves because the expectation passes its total, not because the total does.

There is one thing that scalar must not be asked to do, and we correct it here rather than
leave a row label carrying it. $A$ is a per-target quantity, and $41$ is the \emph{mean} of
the per-target lower bounds $A_j$, which run over $[1,132]$ with quartiles $8$, $33$ and
$66$. Because $\mathbb{E}[K_j] = m_j/(A_j{+}1)$ is convex in $A_j$, substituting that mean
into it understates the null expectation instead of bounding it. Run per target, the honest
lower bound expects $349.0$ exceedances and not $88.2$, with $13$ targets at $A_j{\leq}5$
supplying $228$ of the total between them. So it is the per-target row of
Table~\ref{tab:nullcontrol}, and not the $A{=}41$ row, that a reader should read as the least
favourable admissible budget, and it is the more adverse of the two in exactly the arms whose
margin the scalar flattered. There the exact-match arm stands at $1.1$ times its null rather
than fourfold above it, and its median $p$ falls from $1.000$ to $0.760$; the shared-NLI arm
stays below $0.001$. \textbf{The adjudicator's median $p$ is $1.000$ there too}, on every one
of the $101$ tie draws, with its total of $652$ still standing at $1.9$ times the honest
null. The correction spends margin in the two arms the rule is not stated on and spends none
of the verdict, which is the whole point of having stated the rule on the arm whose answer
does not move with a budget we cannot measure.

\textbf{The decision rule was pre-registered on the adjudicator arm alone} (below), and that
arm's verdict is invariant to the one parameter we cannot pin down. That invariance, and not
any headroom in the shared-NLI arm, is what the non-rejection rests on. Reporting the NLI
arm's $0.939$ beside it is reporting a $p$-value at the end of a range most favourable to the
conclusion we draw, and we say so rather than quote it as though the budget behind it had
been measured.

\paragraph{Decision rule (single, pre-registered, with disclosed deviations).}
Reframe~(a) is the claim that ``the attack targets the detector''. It is claimed iff, at
$n{\geq}80$ with $m{=}50$ benign paraphrases per target, the adjudicator clusterer's
randomised-tie exceedance test rejects at the $5\%$ level in the direction of the attack. That design was
fixed in advance at a simulated power of $0.84$ against a two-fold effect and $0.99$ against
a three-fold one, chosen over a cheaper $m{=}30$ (power $0.67$ at the design-time
simulation settings, $0.66$ on the higher-precision re-run quoted above) because our own
selection-on-noise measurement indicates the true effect is modest and the design must be
able to see a small one. We disclose that those pre-registration figures were the
oracle-calibrated ones; the deployed analytic-null test, cutting at its stricter achieved
level, delivers $0.77$ at $m{=}50$ and $0.51$ at $m{=}30$ for that design. That shortfall is a level rather
than a weaker statistic, since re-cut to the oracle's level the two agree to three decimals,
as above. It is nonetheless the operating point we actually ship. The ordering that drove the choice
is unchanged and we do not re-choose $m$ after the fact, but the design rejects less
readily than the number we picked it on, and
a non-rejection must be read against neither of those figures. Both are computed at $n{=}80$
with a uniform $m$, and the run delivered $n{=}77$ with six short arms; at the realized arm
lengths the deployed power against a two-fold effect is $0.67$, and that is the number a
non-rejection here has to be read against. We state the awkward part rather than let a
reader find it: $0.67$ is also, to two decimals, the oracle power at the design-time settings
of the cheaper $m{=}30$ design this paragraph says was passed over as too weak. The two
quantities are not the same one, the first being the deployed power of the design that ran
and the second the oracle power of a design that did not, but the coincidence is fair to
record, and the honest summary is that the run bought the sensitivity of the option it
declined. If the test does not reject, (a) is not claimed, and the
NLI/exact-match bracket together with the judge's point estimate is reported as the honest
result. We are precise about what is bounded there, because the pre-registration was not: the
bracket bounds the mean intended move, and the judge's point estimate lies outside it on the
paired nets of Table~\ref{tab:nullcontrol} that the closing paragraphs lead with, so the word
\emph{bounded} is not available for the quantities this section concludes from. That is the branch the run took, and
the paragraphs after the deviations report it.

Five deviations from the original pre-registration are disclosed as outcome-triggered or,
in the last case, as an omission, rather than folded silently into ``the pre-registered
rule''. The adjudicator moved from the
sentence embedder to the LLM-judge, triggered by the embedder's $0.51$ hard-negative AUROC.
The benign floor moved from individual draws to a budget-corrected comparison, triggered by
the winner's-curse mismatch surfaced in adversarial review. The budget correction then moved
from brute-force matching to the analytic exceedance null, triggered by the matched design
pricing out at what we then costed as $228$ GPU-hours. We correct that figure here rather than
leave it standing: re-derived from the wall clock of runs that have since executed, the matched
design is about $99$ GPU-hours. The deviation still holds, because an analytic null that costs
nothing is preferable to a four-day matched arm. Even so, the trigger we recorded was
$2.3\times$ too large, and a reader checking whether the rule was tuned is entitled to the
corrected number. And the tie rule moved from
conservative to randomised
breaking, triggered by the simulation above showing the conservative rule has no power under
the ceiling.

The fifth is a deviation by omission, and it is the one a reader auditing the
pre-registration is most likely to catch us on, so we state it rather than let four disclosed
deviations imply that the primary statistic was used as written. The pre-registration made
this test primary \emph{conditionally}. The condition was a null-objective ablation, the same
beam search run against a null objective, and the exact test was to be primary if and only if
that ablation's mean exceedance count came back within $[0.5\times, 2.0\times]$ of the
theoretical $m/(A{+}1)$; below the band the analytic null would be anti-conservative and we
would fall back to a prefix statistic or to a null calibrated from the ablation's own
distribution, and above it the test could stand as conservative with that noted. \textbf{The
ablation was never run.} Its only artifact is a single-target checkpoint written under a
record schema this campaign has since superseded, and its own plan document records its
status as not run. We used the exact test as primary anyway. No trigger licenses that, because
there was no outcome to trigger on: the gate was not failed, it was not evaluated, and the
distinction is one a pre-registration audit is entitled to be told rather than left to
discover. What it was guarding is stated below with the rest of the null's specification.

Each of the first four triggers is objective, each was recorded before the corresponding data
was seen, and the full sequence is given here rather than in an appendix precisely because a
reader is entitled to check that the rule was not tuned to the result.

\paragraph{Result: what the rule returns.}
The confirmatory run completed on 2026-08-29. It attacks the \emph{false-alarm stratum of the
attacked pool}, all $80$ score-independently selected correct-answer targets, which is the
pre-registered $n{\geq}80$ met by this population and by no other quantity in this section. The
run scores each target's outcome under three clustering relations, against $m{=}50$ benign feasible
paraphrases drawn from the same proposer and passed through the same feasibility gate, plus a
three-seed same-question band. The judge is Qwen2.5-7B-Instruct. \textbf{The rule returns a
non-rejection, and it does so in every arm.} In the adjudicator arm the randomised-tie
exceedance test gives a median $p$ of $1.000$ across $101$ tie-break realisations, with no
realisation anywhere near the $5\%$ cut and no admissible budget that moves it; the shared-NLI
and exact-match arms give $0.939$ and $1.000$ at the declared $A{=}181$.
Reframe~(a) is therefore \textbf{not claimed}. The direction is worth stating, but it has to
be stated arm by arm rather than for the design as a whole, because the expectation it is
measured against moves with the budget and the budget is not identified.
The observed total is itself a tie-break realisation, and we quote it as one rather than
setting it beside a median $p$ as though it were fixed. At the seed-$0$ draw it is $30$,
$383$ and $636$ against an expectation of $20.4$ at $A{=}181$. The medians over the same
$101$ draws that the $p$-values summarise are
$28$, $383$ and $652$. The exact-match and adjudicator totals sit \emph{above} their null
expectation at every admissible budget, which is the direction \emph{against} the attack and
is the invariance the verdict rests on. The shared-NLI total does not: it is above
expectation at $A{=}181$, but the expectation catches it near $A{=}131$ and passes it below
that, so over most of the defensible range the NLI statistic sits on the attack's side of its
own null, and below $A{\approx}90$ a majority of tie draws would reject. That is a statement
about a parameter we cannot measure rather than a result, and it is the reason the rule is
stated on the arm that does not depend on it. The shared-NLI total is also the one that moves
across tie draws, over $[16,41]$, and $0.229$ is its lowest $p$ at $A{=}181$ and not in
general. In
the arm the optimiser actually searched, its search behaves like
$n_{\mathrm{eff}}{\approx}122$ random feasible paraphrases. We do not read that as fewer than
the budget it spent: $n_{\mathrm{eff}}$ is denominated in feasible paraphrases while the
$181$ counts scored evaluations, and against the feasible budget's own point estimate of
$121$ it falls short of nothing. Off that arm the same effect size collapses:
$n_{\mathrm{eff}}$ is $8.7$ under exact-match and $4.8$ under the pre-registered adjudicator,
so a search costing $181$ evaluations ranks where the best of about five random benign
paraphrases would. That is not a clean measure of the search's strength, since off the
objective arm the attack's query is re-scored rather than selected and so inherits the
mis-specification set out below, but it is the effect size the pre-registered arm returns and
it points the same way as that arm's $p$.
Doubling the measured tie multiplicity, the one perturbation that moves $p$ downward, leaves
the NLI arm at $0.229$ and the other two at $1.000$.

A non-rejection is not a measurement of zero, and we do not report it as one. At the deployed
operating point the design that ran, $n{=}77$ with six short arms, has power $0.67$ against a
two-fold effect, so the
correct reading is that an effect of that size is \emph{not detectable at this $n$ on this
population}, not that none exists. The same discipline governs every supplementary interval in
this section, and we state it once as a rule so that it can be checked against every later
sentence: an interval that covers zero is reported as \emph{covering zero}, never as null. The
exact-match net covers zero. So does the index-matched adjudicator net of the last paragraph
below, and it gets the same wording rather than a stronger one.

\paragraph{Denominators, and the convention, both disclosed.}
Four different $n$ appear above, all of them inside the false-alarm stratum of the attacked
pool, and none of them may be substituted for another. The campaign
is $80$; every benign-referenced statistic runs
on $77$, because three targets have no benign arm at all; convention~(c) runs on $69$; and the
full $m{=}50$ budget was achieved on $71$. Three of those four are row denominators in
Table~\ref{tab:nullcontrol}, and the table is read block by block: the bands run over all $80$
attacked targets, every benign-referenced row over the $77$ with a non-empty benign arm, and
the convention~(c) row of Table~\ref{tab:nullconventions} over $69$.
The exceedance test and the paired net are both
benign-referenced, so both run on $77$. The three empty arms
(\texttt{qb\_565}, \texttt{qz\_3393}, \texttt{qb\_2689}) are the targets whose every one of
$250$ feasibility attempts was rejected, and on each of them the optimiser also found no
feasible paraphrase. Taken together, those two failures are evidence that no feasible
paraphrase of those
questions exists, rather than that the sampler was unlucky. That second property is not
peculiar to those three: $11$ of the $80$ targets have no feasible candidate at the
optimiser's best, identically identified by three independent criteria (the returned query
equals the original, the optimiser records no improvement, and its feasible-at-best count is
zero). What singles out these three is that the benign sampler failed on them too, and it is
the empty benign arm, not the optimiser's no-op, that removes a target from the denominator.
They are excluded and stated, never absorbed into a denominator of $80$.
Convention~(c) of Table~\ref{tab:nullconventions} drops eight rows, not eleven, and the
arithmetic is worth spelling out because $77{-}11$ is not $69$: three of the $11$ no-op targets
are the three with an empty benign arm, and those three are already outside the $77$, so only
the remaining eight are excluded there. Those no-op targets bear on the paired net for a reason
that is structural rather than statistical. The optimiser returns the original question when
nothing beats it, so the attack's move is exactly $0$ on the $11$ and, in the NLI arm it was
optimised against, is bounded below by $0$ everywhere (measured minimum $+0.0000$ over all
$80$). The feasibility gate forbids the benign arm from drawing the no-op, so its maximum is
unconstrained and is negative on $10$ of $77$ targets under NLI. The paired net therefore
credits the attack with the null's losses unless a convention intervenes.
Six further arms are
short, at $m$ of $1$, $26$, $27$, $27$, $35$ and $38$. The resulting inflation of the
net is $+0.006$ nats as-is and $+0.001$ under the locked convention, because the benign score
distributions are lumpy: the discrete estimator has only $39$ attainable values at $N{=}10$,
and $29\%$ of a target's benign NLI draws already sit at that target's own maximum, so a maximum over
$26$ draws almost always equals the maximum over $50$. That is a disclosure, not a correction.
The two disclosures pull the same way and neither is large enough to move a conclusion: the
short arms inflate the net, and the locked convention deflates it. Finally, and non-negotiably:
convention~(d) of Table~\ref{tab:nullconventions} was fixed in a pre-registration written when
$68$ of the $80$ targets were visible, which locked it in advance of the last $12$ and after
the effect of every candidate convention on
every arm had been computed. It is not a blind pre-registration and must not be read as one.
What can be said for it is stated in that document and repeated here: it is the convention least
favourable to our own attack in all three arms, which is the direction that makes a post-hoc
choice non-self-serving, and all four are reported, as that document requires, so that a reader
can price the choice rather than take it on trust. Of the alternatives, (a) is indefensible on
principle and (c) changes the denominator.

\paragraph{Where the exceedance null is, and is not, correctly specified.}
The attack's objective is the shared-NLI entropy and nothing else, so the query it reports is an
argmax over that target's feasible \emph{NLI} candidates, at most $181$ of them, and the
exact-match and judge values are
independent re-scorings of that single query. The exceedance null prices in a maximum over $A$
exchangeable draws. That is exactly what the NLI arm contains and exactly what the other two do
not: there the attack contributes one unselected observation, so the null expects it to be far
more extreme than it has any reason to be, and $p$ is inflated. The mis-specification runs
against the attack, which is the conservative direction for the claim we are not making, but it
has a consequence a reader should not have to derive: the exact-match and judge non-rejections
are weaker evidence than the NLI one, and the NLI arm is the only place where this test is
asking the question it was designed to ask. That holds even though the NLI arm is itself
confounded as an effect estimate.

A further mis-specification is structural, lives in the code rather than in the data, and is
the one the missing ablation existed to price. The null treats a target's attack candidates
and its benign draws as exchangeable draws from one distribution, but the optimiser proposes
each child from its \emph{beam parent}, so only its first round is single-hop and everything
after it is a second- or higher-order rewrite, while the benign sampler proposes every draw
from the \emph{original question} and so is single-hop throughout. Two effects run against
each other and we do not know the net sign: multi-hop drift would make the attack's later
candidates more dispersed for reasons that have nothing to do with optimisation, which is the
anti-conservative direction, and beam clustering would have the search revisit a small region
and supply fewer effective draws, which is the conservative one. Measuring that net sign is
exactly what the ablation would have done. What can be said without it is that the arm the
rule is stated on sits far from its rejection region under either sign: the adjudicator's
total stands above its null expectation at every budget in $[41,181]$ and at the per-target
lower bound too, at worst by a factor of $1.9$. That bounds the consequence of the asymmetry;
it does not measure it, and we do not offer it as though it did.

That is not the same as resting the verdict on the weakest arm, and a reader who joins this
paragraph to the budget argument above will otherwise conclude that it is. The verdict is the
withholding of a claim rather than a claim, so a mis-specification that inflates $p$ can leave
reframe~(a) unclaimed but can never manufacture it; what the adjudicator's invariance across
$[41,181]$ buys is procedural, that the pre-registered rule returns one answer at every budget
we cannot rule out, while the evidence \emph{about} the attack lives in the NLI arm, where the
honest report is that the $p$-value is not identified.

For one of the three arms that weakness is total, and we state it rather than let three
reported $p$-values read as three confirmations. The deployed cut is $S{\leq}11$: exact
convolution of the per-target nulls gives $\Pr(S{\leq}11) = 0.0296$ and
$\Pr(S{\leq}12) = 0.0501$, so $11$ is the largest total this test can reject on at $n{=}77$.
The exact-match arm cannot produce one. Its benign draws are pinned at the $\log N$ ceiling
on $1{,}461$ of $3{,}704$ draws ($39.4\%$), against $10.7\%$ of the NLI arm's draws and
$1.2\%$ of the judge's; that is a draw-level rate, not the target-level clean saturation
Table~\ref{tab:nullcontrol} reports. Every one of those ties hands the benign side tie
credit. Simulating the strongest
alternative this design exists to catch, an attack that reaches the ceiling on every target at
that target's measured tie multiplicity, the exact-match total has expectation $150$ and its
smallest value over $20{,}000$ draws is about $110$. Its rejection region is unreachable.
Power against that alternative is $0.56$ in the shared-NLI arm and $1.00$ in the
adjudicator's, and $0.00$ in exact-match. The exact-match $p{=}1.000$ is therefore a certainty
of the arithmetic rather than evidence about the attack: the non-rejection rests on two arms
that could have rejected and did not, not on three.

The shared-NLI arm carries a mis-specification of its own, and it runs the other way, so we disclose it
here rather than leave it to the pre-registration. Exchangeability requires the attack's
reported value to be a draw from the same feasible-paraphrase distribution as the benign
draws. On a no-op target it is not: the optimiser returned the original question, so the
attack's move is a structural $0$, while the feasibility gate rejects an identical
``paraphrase'' as its first test, so the benign arm can never produce one. The eight no-op
targets that have a benign arm therefore donate an observed exceedance count of exactly $0$
against an expectation of $1.93$, and they do so in all $101$ tie-break realisations rather than
on average. That pushes $S$ down and $p$ toward rejection \emph{in the attack's favour}. It is
an anti-conservative violation of the one arm the test is correctly budgeted for, it is
unresolved, and at a borderline result it would matter. It does not change this one: with
$1.93$ of the shared-NLI arm's $20.4$ expected exceedances handed to the attack for free, the
test still does not reject at any realisation. The other two arms show the reverse, since their
no-op targets donate far \emph{more} exceedances than expected. Nothing about this
violation is therefore shared across the three.

Nor does the opposite correction rescue the attack. Pricing those arms at a budget of one
would over-correct: the attack's judge value is not an independent draw but the judge score of
an NLI argmax, so its effective budget is whatever the argmax inherits through the correlation
between the two clusterers, somewhere strictly between $1$ and $A$. We do not model that
quantity. The index-matched row of Table~\ref{tab:nullcontrol} sidesteps the need to model it,
by holding the selection rule fixed across both sides, and it covers zero.

\paragraph{What the judge arm does and does not say.}
The paired net is \textbf{supplementary} to the exceedance test. Table~\ref{tab:nullcontrol}
reports it under the pre-registered convention~(d) (Table~\ref{tab:nullconventions}) and
against three comparators, because its sign depends on which one is used. Under \emph{every}
convention the judge column is measured against a benign maximum selected on the judge while
the attack's judge score is unselected; that asymmetry is what puts the judge column below
zero, and the convention is not.
Under the locked convention the paired net in the adjudicator arm is
$-0.173$ [$-0.302$, $-0.039$] nats, an interval lying entirely below zero, and we state the
sign plainly: measured this way, the optimised attack moves the adjudicator's semantic entropy
\emph{less} than the benign comparator does, by about $0.17$ nats per target. The result is
stable to the convention and to the bootstrap. It is negative under (b), (b$'$), (c) and (d),
and it agrees with a $t$-test ($p{=}0.013$) and a signed-rank test ($p{=}0.007$).

It would nonetheless be an error to report that as ``the attack loses to budget-matched random
paraphrasing'', and we do not. The comparator in that row is the benign candidate that maximises
the \emph{judge} score, chosen from $50$; the attack's judge score is an unselected re-scoring of
a query chosen on NLI. The benign arm is handed a search the attack never ran, and that search
is worth
$+0.49$ nats of selection premium in that arm. The NLI-selected benign candidate attains the
judge maximum on only $37$ of those $77$ targets. Scoring both sides the same way, by taking the
benign candidate that maximises NLI and reading \emph{its} judge move, the net is
$+0.016$ [$-0.112$, $+0.150$] with a sign test at $p{=}0.66$; the exact-match arm moves the same
way, from $-0.027$ to $+0.018$. Against the mean benign draw the same data give
$+0.175$ [$+0.045$, $+0.312$], and the attack's judge score sits at the $0.70$ mid-rank
percentile of its target's benign draws. The sign of the adjudicator's ``attack effect'' is
therefore a property of how the null is budgeted and selected, not of the attack: the interval
clears zero above against a single draw, contains zero against a matched selection, and clears
zero below against a judge-targeted best-of-$50$. We do not read the first of those as a win
either. Its sign test does not reject ($p{=}0.44$), so its mean is carried by a minority of
large movers. The null control exists to correct one defect, a maximum compared against a
differently-budgeted maximum, and that defect reappears here one level
down, per arm rather than per design. This time it runs \emph{against} the attack rather than
for it, which is very likely why review aimed at pro-attack bias did not surface it. \textbf{The
defensible reading of the adjudicator arm is that its index-matched interval covers zero}. By
the rule stated above we report it as covering zero rather than as null. An interval
spanning $[-0.112, +0.150]$, about a quarter of a nat, does not license that stronger word.
What it does support is narrower: at this $n$, on this
population,
and with the comparator matched, we cannot distinguish the optimised attack from a
budget-matched benign
search. Even index-matched the attack retains $181$ NLI evaluations against the benign arm's
$50$, so that reading is not a claim in the attack's favour either.

\paragraph{Transfer and defense.}
These results are pending. Both the SE$\leftrightarrow$SRE transfer table and the
defense curve for reformulation-averaging fill from the same campaign. Both are reported net
of the benign floor, and both read as corroborating rather than independent evidence, because
the two detectors share the DeBERTa backbone.
